\section{Доступ к прошлому: наблюдение, сообщение и память}
\label{sec:observer-theory}

Пусть $U\ne\varnothing$ --- множество допустимых историй одного устройства при
фиксированных границах системы и протоколе. Целевой бит $Y(u)\in\{0,1\}$
описывает прошлое устройства в истории $u$. Через $H(u)$ обозначим сведения,
доступные получателю к моменту решения, через $W(u)$ --- сведения отправителя.
Пришедшее сообщение обозначим $Z(u)$. В $H$ можно включать только те прежние
показания, которые доступны получателю по его интерфейсу или сохранены в его
памяти. Два протокола с разным доступом к журналу представляют разные задачи.

\subsection{T1. Критерий точного декодирования}
\label{thm:observer-t1}

Для произвольной карты наблюдения $O:U\to\mathcal O$ точное детерминированное
восстановление $Y$ на всех допустимых историях возможно тогда и только тогда,
когда
\begin{equation}
 \forall u,v\in U\qquad O(u)=O(v)\ \Longrightarrow\ Y(u)=Y(v).
 \label{eq:observer-fibres}
\end{equation}
\emph{Доказательство.} Если $Y=d\circ O$, применение одной функции $d$ к равным
аргументам даёт равные ответы. Обратно, для каждого $o\in O(U)$ множество
$\{Y(u):O(u)=o\}$ непусто и по условию одноэлементно. Определим $d(o)$ его
единственным элементом. Значения $d$ вне $O(U)$ произвольны. Это доказывает
существование функции. На конечной области её можно реализовать таблицей при
наличии памяти и адресации соответствующей ёмкости.

Критерий применим отдельно к $O=H$ и $O=(H,Z)$. Чтобы установить выигрыш от
сообщения в смысле безошибочного доступа, достаточно предъявить пару с общим
$H$ и разными $Y$, а затем дать точный декодер по $(H,Z)$ для \emph{всего}
класса историй. При заданном семействе наблюдений $O_t$ момент первого такого
доступа есть наименьшее $t$, при котором выполняется
\eqref{eq:observer-fibres}, если оно существует. В системе со стиранием
выполнение критерия в этот момент оставляет открытым его выполнение позднее.

\subsection{T2. Цена одного сообщения}
\label{thm:observer-t2}

Пусть $U$ конечно. Отправитель посылает одно слово $c(W)$ по безошибочному
каналу фиксированной длины. Моменты отправки и получения заданы заранее.
Построим неориентированный граф $G$ с вершинами $W(U)$. Между $w,w'$ проведём
ребро, если существуют $u,v\in U$ такие, что
\[
 W(u)=w,\qquad W(v)=w',\qquad H(u)=H(v),\qquad Y(u)\ne Y(v).
\]
Допускается $w=w'$. При наличии петли никакого достаточного сообщения такого
вида нет. В отсутствие петель минимальный размер алфавита и минимальное
число передаваемых двоичных разрядов равны соответственно
\[
 M_{\min}=\chi(G),\qquad b_{\min}=\lceil\log_2\chi(G)\rceil .
\]
\emph{Доказательство.} Для двух концов ребра совпадение сообщений породило бы
одинаковую пару $(H,Z)$ с различными целями, что противоречит T1. Следовательно,
значения всякого точного кода дают правильную раскраску. Петля потребовала бы
разных цветов у одной вершины. Обратно, при любой правильной раскраске две
истории с одинаковым $H$ и разными целями получают разные цвета. Если
отправитель передаёт цвет, T1 строит точный декодер. Для $M$ цветов требуется
и достаточно $\lceil\log_2 M\rceil$ битов фиксированной длины, включая
случай $M=1$ и пустого сообщения.

Нижняя граница $\max_{h\in H(U)}|\{Y(u):H(u)=h\}|$ достигается, если
отправитель знает и $H$, и $Y$. Он может локально пронумеровать возможные
цели для каждого $h$. При ограниченном доступе к истории эта оценка бывает
строго меньше $M_{\min}$. Пусть у получателя записана одна из упорядоченных
пар $ab,ac,bc$, отправитель знает одну выбранную букву этой пары, а целевой
бит указывает прежний выбор первой или второй позиции. Каждая пара требует
двух разных цветов для своих букв. Три буквы образуют треугольник, поэтому
здесь нужны три символа сообщения при двух возможных целевых битах.
Связь одношагового безошибочного кодирования со сторонней информацией и
раскраской графа восходит к
\href{https://doi.org/10.1109/TIT.1976.1055607}{работе Витсенхаузена};
функциональную постановку со сторонней информацией рассматривают также
\href{https://doi.org/10.3390/e26040338}{Шарпене, Ле Трёст и Руми}.

\subsection{T3. Ошибки канала и цена адресуемого запроса}
\label{thm:observer-t3}

Пусть в истории $u$ искажённый сигнал может принять любое значение из
непустого множества $\mathcal R(u)$. Восстановление при \emph{каждом}
допустимом искажении возможно ровно тогда, когда
\[
 H(u)=H(v),\quad Y(u)\ne Y(v)
 \quad\Longrightarrow\quad \mathcal R(u)\cap\mathcal R(v)=\varnothing .
\]
Общий символ из пересечения требует двух разных ответов на одно наблюдение.
Если все указанные пересечения пусты, каждому достижимому $(h,r)$
соответствует единственное значение цели. Это доказательство T1 на расширенных
историях с учётом исхода канала.

Для конечного совместного распределения $(H,R,Y)$ и штрафа за неверный бит
минимальная вероятность ошибки имеет другой смысл и равна
\[
 P_e^*=1-\sum_{h,r}\max_{y\in\{0,1\}}\Pr(H=h,R=r,Y=y).
\]
Действительно, детерминированный ответ $d(h,r)$ даёт вероятность успеха
$\sum_{h,r}\Pr(H=h,R=r,Y=d(h,r))$. Каждый член максимизируется независимо.
Например, если $H,Z$ --- равновероятные независимые биты,
$Y=H\oplus Z$ и $R=Z\oplus N$ с независимым $N\sim\mathrm{Bernoulli}(p)$,
$0\le p\le 1/2$, то оптимальный ответ $H\oplus R$ ошибается с вероятностью
$p$. При $p>0$ точная гарантия для всех искажений теряется.

Преимущество интерактивности зависит от того, какие сообщения разрешены.
Пусть $A_0,A_1$ --- два ранее записанных свободных бита, а локальный адрес
получателя $K$ определяет требуемый бит прошлого $Y=A_K$. Все восемь троек
$(K,A_0,A_1)$ допустимы. Достаточно одного бита запроса $K$ и одного
обратного бита $A_K$. Если обратный бит обязан быть фиксированной функцией
$c(A_0,A_1)$ независимо от адреса, то по одному его значению можно было бы
запустить декодер при $K=0$ и при $K=1$ и восстановить сразу пару
$(A_0,A_1)$. Четыре различные пары требуют не менее четырёх сообщений.
Таким образом, отправка адреса уменьшает необходимый обратный алфавит при
заданном устройстве обмена. Два обратных бита уже позволяют передать всю
запись без запроса. В рамках одного неизменного запроса заранее
воспроизведённый верный ответ с тем же содержанием и временем доставки
приводит к тому же исходу у детерминированного получателя.

\subsection{T4. Собственный отчёт и сохранённое состояние}
\label{thm:observer-t4}

Рассмотрим два устройства, которые вычисляют один и тот же точный
$d(H,Z)=Y$ и публикуют его. Первое сохраняет бит $M=Y$. Второе стирает
свидетельство и попадает в одно и то же доступное состояние хотя бы для
двух историй с разными $Y$. Затем прежний публичный отчёт и канал наблюдателя
закрываются, обоим устройствам предъявляется новый бит $C$ и запрашивается
действие $Y\oplus C$. Первое устройство вычисляет $M\oplus C$. Для второго
зафиксируем одно и то же $C$ на указанной паре историй. Полные доступные
входы равны, а требуемые биты противоположны. По T1 у второго устройства
точного декодера нет. Так две реализации с одинаковым немедленным отчётом
различаются отложенным действием.

Вмешательство в $Z$ при фиксированном раннем состоянии проверяет зависимость
действия от доставленного сигнала. Для проверки переноса результата через
эпизод требуется второе вмешательство: подмена или стирание $M$ при одном
и том же новом запросе и исключённом доступе к прежнему отчёту. Это два
разных контрфактических сравнения. На приведённой задаче одного внутреннего
регистра достаточно для успешного отложенного действия.
